\chapter[Transformers, VLMs, and Grounding]{Transformers, VLMs,\\and Grounding}

\section*{학습 목표}

이 장은 Transformer와 Vision-Language Model(VLM)이 VLA의 backbone으로 들어오는 경로를 설명한다. 목적은 Transformer 일반론을 길게 반복하는 것이 아니다. Real VLA 관점에서 필요한 질문은 더 구체적이다. Tokenization은 robot action representation에 어떤 제약을 주는가? VLM의 image-text grounding은 robot manipulation grounding과 어떻게 다른가? Web-scale semantic prior는 어떤 부분을 도와주고, 어떤 부분은 여전히 robot data와 controller가 맡아야 하는가?

\begin{enumerate}
  \item tokenization과 self-attention을 VLA에 필요한 수준으로 정리한다.
  \item ViT와 VLM의 multimodal representation을 설명한다.
  \item contrastive image-text alignment와 open-vocabulary recognition의 의미를 정식화한다.
  \item semantic grounding, visual grounding, affordance grounding, physical grounding을 구분한다.
\end{enumerate}

\section{Why VLM matters for VLA}

VLA가 가능해진 배경은 로봇공학 내부에만 있지 않다. Transformer는 sequence modeling을 확장했고, ViT는 이미지를 patch token sequence로 바꾸었으며, VLM은 image와 text를 같은 representation space에서 정렬했다 \cite{vaswani2017,dosovitskiy2021vit,clip2021,flamingo2022,blip22023}. 이 흐름은 로봇 policy가 open-vocabulary instruction과 visual scene을 조건으로 action을 생성할 수 있게 만든 핵심 배경이다.

하지만 VLM은 그 자체로 VLA가 아니다. VLM은 보통 다음 함수를 학습한다.
\begin{equation}
  \bm{h}^{vlm}
  =
  F_{\theta}(\bm{I},\bm{x}^{text}),
  \label{eq:vlm-representation}
\end{equation}
여기서 $\bm{I}$는 image, $\bm{x}^{text}$는 text token sequence이다. VLA는 이 representation을 robot action으로 연결해야 한다.
\begin{equation}
  \act_t
  =
  G_{\phi}(\bm{h}^{vlm}_t,\conf_t,\bel_t).
  \label{eq:vlm-to-action}
\end{equation}
식 \eqref{eq:vlm-to-action}의 $G_{\phi}$가 바로 robot-specific policy, action head, planner, controller interface이다. 이 부분이 없으면 VLM은 scene을 설명할 수 있어도 robot을 움직이는 model은 아니다.

\section{Tokenization}

Transformer는 입력을 token sequence로 본다. text는 subword token sequence로 표현된다.
\begin{equation}
  \bm{x}^{text}
  =
  (x_1,x_2,\ldots,x_N),
  \qquad
  x_i\in\mathcal{V}.
  \label{eq:text-token-sequence}
\end{equation}
image도 patch token으로 바꿀 수 있다. image $\bm{I}\in\R^{H\times W\times C}$를 $P\times P$ patch로 나누면 patch 수는
\begin{equation}
  N_p = \frac{HW}{P^2}.
  \label{eq:number-image-patches}
\end{equation}
각 patch를 embedding으로 projection하여
\begin{equation}
  \bm{z}^{img}_i
  =
  E_{img}\mathrm{vec}(\bm{I}_i)+\bm{p}_i
  \label{eq:vit-patch-embedding}
\end{equation}
를 만든다. 여기서 $\bm{p}_i$는 positional embedding이다.

VLA에서 tokenization의 의미는 더 넓다. language token, image patch token, proprioception token, history token, action token이 하나의 sequence에 들어갈 수 있다.
\begin{equation}
  \bm{Z}
  =
  [\bm{z}^{text}_{1:N},\bm{z}^{img}_{1:P},\bm{z}^{prop},\bm{z}^{hist},\bm{z}^{act}_{1:K}].
  \label{eq:vla-token-sequence}
\end{equation}
이 구조는 flexible하지만, robot action이 token으로 잘 표현되는지는 별도 문제이다. Discrete token은 VLM/LLM pipeline과 잘 맞지만, continuous contact-rich manipulation에서는 quantization, latency, smoothness 문제가 생길 수 있다.

\section{Self-attention}

Self-attention은 token 간 관계를 계산한다. token embedding matrix를 $X\in\R^{N\times d}$라 하면,
\begin{equation}
  Q=XW_Q,\qquad K=XW_K,\qquad V=XW_V.
  \label{eq:qkv}
\end{equation}
attention output은
\begin{equation}
  \mathrm{Attn}(X)
  =
  \mathrm{softmax}\left(\frac{QK^\top}{\sqrt{d_k}}\right)V.
  \label{eq:self-attention}
\end{equation}

VLA에서 attention의 장점은 서로 다른 modality 사이의 관계를 모델링할 수 있다는 점이다. 예를 들어 ``빨간 컵을 잡아라''라는 text token은 image patch의 red object region과 연결될 수 있고, ``잡아라''라는 verb는 gripper action과 연결될 수 있다. 그러나 attention weight가 높다고 해서 물리적으로 grasp 가능한 것은 아니다. Attention은 representation relation이고, feasibility는 robot geometry와 contact constraint의 문제이다.

\section{Contrastive image-text alignment}

CLIP-style contrastive learning은 image encoder와 text encoder를 같은 embedding space에 정렬한다. image embedding을 $\bm{u}_i$, text embedding을 $\bm{v}_j$라 하면 similarity는
\begin{equation}
  s_{ij}
  =
  \frac{\bm{u}_i^\top \bm{v}_j}
  {\|\bm{u}_i\|_2\|\bm{v}_j\|_2}.
  \label{eq:clip-similarity}
\end{equation}
batch 안의 paired image-text를 맞추는 contrastive loss는 다음처럼 쓸 수 있다.
\begin{equation}
  \mathcal{L}_{i2t}
  =
  -\frac{1}{B}
  \sum_{i=1}^{B}
  \log
  \frac{\exp(s_{ii}/\tau)}
  {\sum_{j=1}^{B}\exp(s_{ij}/\tau)}.
  \label{eq:clip-loss}
\end{equation}
text-to-image 방향 loss도 대칭적으로 더할 수 있다.

이 objective는 image와 caption 사이의 semantic alignment를 강하게 만든다. 그래서 VLM은 zero-shot object recognition, semantic retrieval, open-vocabulary grounding에서 강하다. 그러나 이 loss는 grasp stability, friction, mass, force limit, controller compatibility를 직접 학습하지 않는다.

\section{Multimodal fusion}

VLM은 image token과 text token을 여러 방식으로 결합할 수 있다. 간단한 late fusion은
\begin{equation}
  \bm{h}
  =
  \mathrm{Fuse}(\phi_I(\bm{I}),\phi_T(\bm{x}^{text}))
  \label{eq:late-fusion}
\end{equation}
로 쓸 수 있다. Interleaved multimodal model은 image/video/text token을 하나의 sequence로 처리한다.
\begin{equation}
  \bm{Z}
  =
  [\bm{z}^{text}_{1:n_1},\bm{z}^{img}_{1:p_1},\bm{z}^{text}_{n_1+1:n_2},\bm{z}^{img}_{p_1+1:p_2}].
  \label{eq:interleaved-sequence}
\end{equation}

Robot setting에서는 여기에 proprioception과 action history가 추가되어야 한다.
\begin{equation}
  \bm{h}^{robot}_t
  =
  F_{\theta}(\bm{I}_{t-H:t},\bm{x}^{text},\conf_{t-H:t},\act_{t-H:t}).
  \label{eq:robot-multimodal-fusion}
\end{equation}
이 식은 VLA가 단순 VLM fine-tuning이 아니라는 점을 보여준다. Robot policy에는 body state와 action history가 필요하다.

\section{Four kinds of grounding}

Grounding이라는 단어는 자주 너무 넓게 쓰인다. Real VLA에서는 최소 네 가지 grounding을 구분해야 한다.

\begin{table}[ht]
\centering
\caption{Grounding의 네 가지 의미}
\label{tab:grounding-types}
\small
\begin{tabularx}{\textwidth}{p{31mm}X X}
\toprule
종류 & 의미 & VLA에서 필요한 추가 조건 \\
\midrule
Semantic grounding & 단어와 개념을 연결한다 & instruction이 task와 constraint로 변환되어야 한다 \\
Visual grounding & 단어와 image region/object를 연결한다 & object pose, affordance, occlusion을 다뤄야 한다 \\
Affordance grounding & object가 어떤 action을 허용하는지 추정한다 & robot morphology와 current state에 의존한다 \\
Physical/action grounding & action이 실제 물리 결과로 이어진다 & controller, contact, force, safety, feedback loop가 필요하다 \\
\bottomrule
\end{tabularx}
\end{table}

VLM이 주로 강한 영역은 semantic grounding과 visual grounding이다. Robot deployment에서 필요한 것은 affordance grounding과 physical/action grounding이다. 예를 들어 VLM이 ``서랍 손잡이''를 찾을 수 있어도, 그 손잡이가 현재 robot reachability 안에 있는지, 어떤 방향으로 당겨야 하는지, 필요한 force가 얼마인지, gripper가 slip 없이 잡을 수 있는지는 별도 문제이다.

\section{From VLM prior to VLA policy}

VLM prior는 VLA policy에 여러 방식으로 들어간다.

\begin{enumerate}
  \item visual encoder로 사용되어 object와 relation feature를 제공한다.
  \item language encoder로 사용되어 instruction embedding을 만든다.
  \item embodied reasoning layer로 사용되어 subgoal, constraint, plan proposal을 만든다.
  \item action policy backbone으로 fine-tuning되어 action token 또는 action chunk를 생성한다.
\end{enumerate}

이 구조를 하나의 equation으로 쓰면 다음과 같다.
\begin{equation}
  \act_t
  =
  \pi_{\phi}
  (
  F_{\theta}^{vlm}(\bm{I}_t,\task),
  \conf_t,
  \bel_t
  ).
  \label{eq:vlm-prior-policy}
\end{equation}
여기서 $\theta$는 web-scale 또는 multimodal pretraining에서 온 parameter이고, $\phi$는 robot action head 또는 adapter parameter일 수 있다.

VLA 설계의 핵심 질문은 $\theta$와 $\phi$의 경계를 어디에 둘 것인가이다. VLM backbone을 고정하고 action head만 학습할 수도 있고, multimodal backbone 전체를 robot data로 fine-tuning할 수도 있다. 전자는 안정적이지만 embodiment-specific adaptation이 제한될 수 있고, 후자는 강력하지만 데이터와 safety 검증 비용이 커진다.

\section{Grounding gap}

VLM이 잘하는 것과 robot이 해야 하는 것 사이에는 gap이 있다.
\begin{equation}
  \mathrm{GroundingGap}
  =
  \mathrm{Semantics}(\bm{I},\task)
  -
  \mathrm{ExecutableAction}(\state,\conf,\mathcal{C}).
  \label{eq:grounding-gap}
\end{equation}
이 식은 엄밀한 metric이 아니라 책임 경계를 표현하는 notation이다. 첫 항은 VLM이 강한 semantic interpretation이고, 두 번째 항은 robot state, constraint, controller를 만족하는 executable action이다.

이 gap은 다음 failure로 나타난다.

\begin{itemize}
  \item object는 맞게 찾았지만 graspable pose를 잘못 추정한다.
  \item instruction은 이해했지만 forbidden region을 hard constraint로 처리하지 않는다.
  \item open-vocabulary object는 인식하지만 robot reachability를 모른다.
  \item action token은 생성하지만 controller가 고주파로 추적할 수 없다.
  \item caption-level relation은 맞지만 contact force나 hinge direction을 모른다.
\end{itemize}

\section{Evaluation metadata}

VLM-to-robot grounding을 평가하려면 일반 VLM benchmark와 다른 metadata가 필요하다.

\begin{itemize}
  \item object grounding accuracy and pose error
  \item relation grounding: left/right/on/in/behind/inside 같은 spatial relation
  \item affordance prediction: graspable, pushable, openable, reachable
  \item executable action success: controller tracking, contact success, recovery
  \item semantic ambiguity handling: ask, abstain, clarify, safe stop
  \item distribution shift: object category, viewpoint, lighting, layout, distractor
  \item physical metadata: robot morphology, camera calibration, controller, force limit
\end{itemize}

이 metadata가 없으면 VLM의 높은 recognition score를 robot capability로 착각하기 쉽다.

\begin{casebox}{Case study: semantic grounding과 action grounding의 차이}
VLM이 ``red mug''를 정확히 찾고 ``put it in the rack''을 이해해도, mug pose covariance, rack slot reachability, gripper approach direction, forbidden glass region을 action constraint로 변환하지 못하면 robot task는 실패한다. 이 장의 핵심은 VLM grounding이 action grounding의 입력일 뿐이라는 점이다.
\end{casebox}

\chapterwork
{Transformer, ViT, CLIP, Flamingo, BLIP-2를 VLA backbone의 representation source로 읽되, physical grounding과 구분하라 \cite{vaswani2017,dosovitskiy2021vit,clip2021,flamingo2022,blip22023}.}
{attention map이나 image-text score가 높아도 robot affordance claim이 되지 않는 이유는 무엇인가?}
{하나의 instruction에 대해 semantic grounding, visual grounding, affordance grounding, physical/action grounding을 분리한 annotation sheet를 설계하라.}

\section{요약}

Transformer와 VLM은 VLA의 핵심 표현학습 기반이다. Tokenization과 attention은 image, text, proprioception, history, action을 하나의 sequence로 결합할 수 있게 한다. Contrastive VLM은 open-vocabulary semantic grounding을 제공한다. 그러나 semantic grounding은 physical/action grounding이 아니다. Real VLA는 VLM prior를 robot state, controller, contact, safety, evaluation loop와 연결해야 한다.

다음 장에서는 이 VLM/LLM을 robot planner로 사용하던 흐름을 다룬다. 그 시기는 완전한 VLA는 아니었지만, 오늘날 hierarchical VLA와 embodied reasoning architecture의 직접적인 조상이다.
